%%%%%%%%%%%%%%%%%%%%%%%%%%%%%%%%%%%%%%%%%%
% Appendix A — Regime Detection LLM Prompt
% Added in response to Reviewer 3 comment R3.4a:
% "The exact prompts used for the LLM must be provided (preferably in
%  an appendix)."
%%%%%%%%%%%%%%%%%%%%%%%%%%%%%%%%%%%%%%%%%%

\appendix

\section[Regime Detection LLM Prompt]{Regime Detection LLM Prompt}
\label{app:prompt}

This appendix reproduces the complete prompt submitted to the LLM for
each 30-day regime-detection window, together with the API configuration
and output schema, so that the experiment is fully reproducible from the
published manuscript without reference to the source repository.

The authoritative implementation lives at
\texttt{src/llm/mechanics\_prompt\_builder.py::build\_regime\_prompt()}
in the code release accompanying this paper; the text below is
transcribed verbatim from that implementation.

\subsection{Model and API Configuration}
\label{app:prompt:config}

All 2{,}221 evaluations were obtained from OpenAI's \texttt{o4-mini}
reasoning model through the OpenAI Batch API (asynchronous, 24-hour SLA,
100\% completion rate observed across the five validation phases).

\begin{itemize}
    \item \textbf{Model:} \texttt{o4-mini}
    \item \textbf{Temperature:} 1.0. OpenAI reasoning models
          (\texttt{o1}, \texttt{o3}, \texttt{o4} families, and the newer
          GPT-5 reasoning variants) reject user-supplied
          \texttt{temperature} or \texttt{top\_p} values and run at the
          default temperature of 1; our batch submission code does not
          override this, so all 2{,}221 requests used temperature 1
          implicitly.
    \item \textbf{Maximum completion tokens:} not explicitly set in the
          Batch API request body; the OpenAI API default for
          \texttt{o4-mini} applies.
    \item \textbf{Response format:} not enforced via the API
          \texttt{response\_format} field; the model is instructed in
          the prompt to return a JSON object with a specific schema. Of
          the 1{,}307 per-window records for which raw JSON responses
          are retained, 1{,}301 parsed cleanly (99.54\%); six failed
          JSON parsing and are recorded with an explicit \texttt{error}
          field and treated as non-detections.
    \item \textbf{Seed:} the OpenAI Batch API exposes a \texttt{seed}
          parameter for best-effort reproducibility; we did not set a
          seed in this study, so each evaluation reflects the model's
          native sampling at temperature 1.
    \item \textbf{Access mode:} OpenAI Batch API (asynchronous; 24-hour
          SLA per batch submission).
\end{itemize}

\noindent\textbf{Reproducibility note.}
Exact bit-identical replication of any single response is not
guaranteed: temperature is fixed at 1 on the server, and even with a
seed the OpenAI documentation states that determinism is best-effort
and can shift when the server \texttt{system\_fingerprint} changes.
Reproducibility at the \emph{distributional} level is achieved by
(i)~the large sample size (N = 2{,}221 evaluations) and
(ii)~the mechanical criteria embedded in the prompt, which give the
model concrete numerical thresholds to apply rather than asking for
free-form judgment. Section~\ref{sec:regime} reports detection rates
with bootstrap 95\% confidence intervals to quantify the residual
sampling variation.

\subsection{System Message and User Prompt}
\label{app:prompt:body}

The prompt is delivered as a single user-role message (the
\texttt{o4-mini} reasoning model treats the first paragraph as the
de~facto system instruction). The placeholder
\texttt{\{gex\_data\_table\}} is replaced at runtime with the
obfuscated 30-day GEX sequence: one line per day, in the format
\texttt{Day T-29: +3.42B} through \texttt{Day T+0: -12.18B}, where the
calendar date has been replaced with a relative day label and the
ticker symbol is absent. No other identifying context is supplied.

\begin{verbatim}
You are a market structure analyst specializing in dealer gamma
positioning regimes.

TASK: Analyze this 30-day period and determine if it represents a
PERSISTENT regime where dealer constraints create forced, directional
flows.

## 30-DAY GEX DATA

{gex_data_table}

## REGIME CLASSIFICATION FRAMEWORK

### PERSISTENT REGIMES (Detect These)

**1. PERSISTENT POSITIVE REGIME**
- Definition: Dealers are LONG gamma, forced to sell into strength
- Criteria:
  * >70% of days (21+/30) have positive net GEX
  * Average magnitude >$5B
  * <=5 sign flips across 30 days
  * Stable directional constraint

**Mechanism**: When dealers hold long gamma:
- Price rises -> Dealers MUST sell shares (rebalance)
- Price falls -> Dealers MUST buy shares (rebalance)
- Creates dampening, mean-reverting flows
- Constraint is STRUCTURAL (dealers cannot avoid)

**2. PERSISTENT NEGATIVE REGIME**
- Definition: Dealers are SHORT gamma, forced to buy into strength
- Criteria:
  * >70% of days (21+/30) have negative net GEX
  * Average magnitude >$5B
  * <=5 sign flips across 30 days
  * Stable directional constraint

**Mechanism**: When dealers hold short gamma:
- Price rises -> Dealers MUST buy shares (chase)
- Price falls -> Dealers MUST sell shares (chase)
- Creates amplifying, momentum flows
- Constraint is STRUCTURAL (dealers cannot avoid)

---

### NON-REGIMES (Reject These)

**3. TRANSITIONAL (Reject)**
- Frequent sign flips between positive/negative GEX
- No dominant regime direction (less than 70% same sign)
- Market in regime change period
- Example: 15 positive days, 15 negative days (50/50 split)

**Why Reject**: No persistent constraint. Dealers face mixed
conditions daily. Not a structural regime.

**4. LOW CONVICTION (Reject)**
- Consistent sign BUT weak magnitude (<$5B average)
- Example: 25 days positive, avg $2B GEX
- Insufficient constraint to create persistent forced flows

**Why Reject**: Even if sign is consistent, magnitude too weak to
force dealers into meaningful positions. Not a structural constraint.

---

## ANALYSIS QUESTIONS

Systematically evaluate the 30-day window:

**Step 1: Sign Persistence**
1. Count days with positive net GEX
2. Count days with negative net GEX
3. Calculate persistence percentage:
   max(positive_days, negative_days) / 30 * 100
4. Does it meet 70% threshold (21+ days)?

**Step 2: Magnitude Assessment**
1. Calculate average GEX magnitude (absolute value):
   sum(|net_gex|) / 30
2. Is average magnitude >=$5B?
3. Check for extreme outliers that might distort average

**Step 3: Stability Check**
1. Count sign flips: How many times does GEX switch from
   pos->neg or neg->pos?
2. Are there <=5 sign flips across 30 days?
3. Stable regime should have low flip count

**Step 4: Regime Classification**
- If Steps 1, 2, 3 all pass AND positive dominates
    -> PERSISTENT POSITIVE
- If Steps 1, 2, 3 all pass AND negative dominates
    -> PERSISTENT NEGATIVE
- If Step 1 passes but Step 2 fails -> LOW CONVICTION (reject)
- If Step 1 fails -> TRANSITIONAL (reject)

---

## CONFIDENCE CALIBRATION (Mechanical Guidance)

Use these concrete anchors to calibrate confidence:

**90-100 (Very High Confidence)**
- 25-30 days same sign (83-100% persistence)
- Average magnitude >$10B
- 0-2 sign flips (highly stable)
- Example: "29 negative days, avg $15B, 1 flip"

**70-89 (High Confidence)**
- 21-24 days same sign (70-80% persistence)
- Average magnitude $5-10B
- 2-4 sign flips (moderately stable)
- Example: "23 negative days, avg $7B, 3 flips"

**50-69 (Borderline - Use with Caution)**
- 18-20 days same sign (60-67% persistence)
- Average magnitude $3-5B
- 5-7 sign flips
- Example: "20 negative days, avg $4B, 6 flips"
- **Note**: Borderline cases should generally be REJECTED unless other
  factors strengthen confidence

**0-49 (Reject - Not Persistent)**
- <18 days same sign (<60% persistence)
- OR average magnitude <$3B
- OR >7 sign flips
- These are NOT persistent regimes

**Important**: Confidence is a FILTER, not a probability. Use it to
distinguish clear regimes (70+) from borderline (50-69) from noise
(<50).

---

## OUTPUT FORMAT (JSON)

Provide your analysis in this exact JSON structure:

```json
{
    "regime_detected": true/false,
    "regime_type": "persistent_positive|persistent_negative|transitional|low_conviction",
    "positive_days": <count as integer>,
    "negative_days": <count as integer>,
    "avg_magnitude_billions": <value as number>,
    "sign_flips": <count as integer>,
    "persistence_pct": <percentage as number>,
    "confidence": <integer 0-100>,
    "reasoning": "Explain step-by-step why this is/isn't a persistent regime. Reference specific metrics (persistence %, avg magnitude, sign flips). If rejecting, state which criterion failed."
}
```

**IMPORTANT**: All numeric fields (confidence, positive_days,
negative_days, sign_flips, avg_magnitude_billions, persistence_pct)
MUST be numbers (integers or decimals), NOT words like "thirty-five"
or "fifty".

**regime_detected Rules**:
- `true` ONLY if regime_type is "persistent_positive" or
  "persistent_negative"
- `false` if regime_type is "transitional" or "low_conviction"

---

## KEY PRINCIPLES

1. **Selectivity is Expected**: Most windows will NOT be persistent
   regimes (expect 30-50% detection rate)

2. **ALL Criteria Must Pass**: Persistence + Magnitude + Stability
   required for detection

3. **Rejection is Valid**: Saying "no persistent regime" is a correct
   answer for transitional/weak periods

4. **Mechanical Over Qualitative**: Use concrete thresholds
   (70%, $5B, 5 flips) rather than subjective judgment

5. **Structural Focus**: Only detect when dealers are FORCED into
   directional positions by constraints

Analyze the 30-day GEX data above and provide your regime
classification in JSON format.
\end{verbatim}

\subsection{Output Schema and Parsing}
\label{app:prompt:output}

Each response is parsed into the following fields (types shown in
parentheses):

\begin{itemize}
    \item \texttt{regime\_detected} (boolean) --- \texttt{true} only when
          \texttt{regime\_type} is \texttt{persistent\_positive} or
          \texttt{persistent\_negative}; \texttt{false} otherwise.
    \item \texttt{regime\_type} (string) --- one of
          \texttt{persistent\_positive}, \texttt{persistent\_negative},
          \texttt{transitional}, \texttt{low\_conviction}.
    \item \texttt{positive\_days}, \texttt{negative\_days},
          \texttt{sign\_flips} (integers in [0, 30]).
    \item \texttt{avg\_magnitude\_billions} (float, USD billions).
    \item \texttt{persistence\_pct} (float, percentage).
    \item \texttt{confidence} (integer 0--100).
    \item \texttt{reasoning} (string) --- free-form step-by-step
          explanation; retained for the post-hoc reasoning-quality audit
          reported in Section~\ref{sec:regime}.
\end{itemize}

\noindent
Parsing is performed by \texttt{src/validation/batch\_regime\_validator.py}
via a robust JSON extractor that tolerates markdown code-fence wrappers
and minor formatting drift. Across the 1{,}307 per-window records for
which the raw responses are retained in the results YAML files (Phases
1--4 and the Phase 2 negative-control suite), the JSON parse-failure
rate was 0.46\% (6~windows failed to parse as valid JSON and are
recorded with an explicit \texttt{error} field; these windows are
treated as non-detections in all aggregate rates reported in
Section~\ref{sec:regime}). Phase 5 multi-year per-window records were
not retained in the published pipeline; Table~\ref{tab:phase5} reports
only the aggregate count per year.
